\section{母池、筛选与证据层级}

全市场研究采用五层漏斗。第一层是申万31个一级行业；第二层是原始预告表949条记录；按A股范围排除4只B股的12条指标后，保留937条合格输入；第三层将归母预告透视为364家A股公司并完成31行业映射；第四层按统一高影响规则筛出73只；第五层结合一年位置、20/60日趋势、扣非纯度、Q1财务和券商时效形成16只优先证据池。

\begin{exhibitbox}[Exhibit 4：全市场研究漏斗]
\begin{tabularx}{\textwidth}{L{2.2cm}R{1.5cm}X}
\toprule
\textbf{层级} & \textbf{数量} & \textbf{规则与用途} \\
\midrule
申万行业 & 31 & 行业位置、资金阶段与预告密度全覆盖 \\
预告公司 & 364 & 以归母净利润预告为公司主键，扣非、EPS和营收作为质量字段 \\
高影响候选 & 73 & 利润规模、同比、扣非纯度与行业阶段统一评分；不先手工挑主题 \\
优先证据池 & 16 & 价格位置、完整历史、正增长、Q1财务与最新研报二次过滤 \\
正式估值池 & 5 & 当前价、股本、2026E分母、三情景、原始PDF Street锚全部可复算 \\
\bottomrule
\end{tabularx}
\end{exhibitbox}

\section{高影响与优先池规则}

高影响定义为：H1归母净利至少20亿元；或至少5亿元、同比至少100\%、扣非占比至少70\%；或至少2亿元、同比至少150\%，且行业属于静默、启动确认或资金观察。高影响只代表值得研究，不代表可买。

进入优先池还需满足正增长和完整一年价格历史。静默行业且位置不高者归“静默吸纳”；一年位置不高、盈利质量过关者归“低位盈利”；行业处于启动／资金观察且20日或60日涨幅达到5\%者归“已启动仍有空间”。负增长、历史不足、扣非不正或一次性收益超过40\%者降级或排除。

\begin{exhibitbox}[Exhibit 5：73只候选最终处置]
\begin{tabularx}{\textwidth}{L{4.5cm}R{1.5cm}X}
\toprule
\textbf{处置} & \textbf{数量} & \textbf{含义} \\
\midrule
静默吸纳优先 & 1 & 行业和位置满足静默条件，仍需估值验证 \\
低位盈利优先 & 11 & 价格位置低且预告质量较好，资金确认可能不足 \\
已启动仍有空间 & 4 & 价格启动且盈利证据较强，需约束追涨与H2兑现 \\
盈利验证观察 & 41 & 高影响但位置、行业阶段或现金流尚未同时通过 \\
价格先行 & 12 & 业绩兑现但一年位置或涨幅偏高，等待回撤 \\
负增长／历史不足／一次性排除 & 4 & 不进入优先池 \\
\bottomrule
\end{tabularx}
\end{exhibitbox}

\section{数据质量与边界}

申万分类使用官方历史分类文件并按2026-07-11之前的最新生效记录映射。本机证书链无法验证上游HTTPS证书，下载过程使用\texttt{verify=False}，但归档文件大小1,161,216字节、SHA-256为\texttt{8da3f757}\allowbreak\texttt{895a6d77}\allowbreak\texttt{a19d8f69}\allowbreak\texttt{0cca3b30}\allowbreak\texttt{22fa0da5}\allowbreak\texttt{6533cdb4}\allowbreak\texttt{769f5939}\allowbreak\texttt{b4bc49d2}。这是传输验证降级，不是内容缺失。4只B股已从A股母池排除；1只北交所A股公司使用明确手工映射并保留。

“主力资金”仍是按成交单大小划分的公开统计，不等于机构身份穿透。浙江东方券商元数据接口返回\texttt{KeyError('infoCode')}，其Q1财务和预告完整保留；进一步检查官方H1预告、Q1财务和公开检索后，仍未取得可核验的当前原始券商PDF，因此Street目标权重为0，House PB-ROE／盈利桥不受影响。16只优先池财务成功16只，当前原始PDF覆盖15只；第16只以正式回退路径闭环。
